\chapter{Firmware}

Firmware finálního prototypu je určen především pro ovládání analogové části, vzorkování napěťového a proudového signálu a následný výpočet impedance. 
Je napsán pro mikrokontrolér STM32F103C8 v prostředí PlatformIO s frameworkem Arduino. 
Pro části, u kterých je důležité přesné časování, je však použita přímo knihovna STM32 HAL.

Program po spuštění inicializuje komunikační rozhraní, řídicí výstupy, zdroj měřicího signálu, indikační LED a snímací část. 
Poté  v hlavní smyčce pouze zpracovává příkazy ze sériové linky. 
Samotné měření tedy není spuštěno trvale, ale vždy až na požadavek z počítače.

\section{Komunikační rozhraní}

Zařízení je ovládáno textovým příkazovým rozhraním po sériové lince. 
Příkazy jsou ukončeny znakem konce řádku. 
Rozhraní slouží jak pro nastavování generátoru a analogového frontendu, tak pro vyvolání snímání nebo výpočtu impedance.

Mezi hlavní příkazy patří:
\begin{table}[H]
    \centering
    \small
    \begin{tabular}{ll}
        \toprule
        Příkaz & Význam \\
        \midrule
        \texttt{status} & Vypíše aktuální nastavení. \\
        \texttt{dds <frekvence> <povoleni>} & Nastaví frekvenci a zapnutí zdroje měřicího signálu. \\
        \texttt{amp <hodnota>} & Nastaví amplitudu měřícího signálu. \\
        \texttt{offset <hodnota>} & Nastaví stejnosměrnou složku signálu pomocí PWM. \\
        \texttt{range <hodnota>} & Vybere bočník pro změnu měřícího rozsahu. \\
        \texttt{vpga <hodnota>} & Nastaví zesílení napěťového kanálu. \\
        \texttt{ipga <hodnota>} & Nastaví zesílení proudového kanálu. \\
        \texttt{capture} & Sejme a odešle jeden z měřicích kanálů. \\
        \texttt{capturepair} & Sejme a odešle oba měřicí kanály. \\
        \texttt{measure} & Sejme oba kanály a vypočítá impedanci. \\
        \bottomrule
    \end{tabular}
    \caption{Přehled podporovaných příkazů firmwaru}
    \label{tab:firmware_prikazy}
\end{table}

Odpovědi jsou záměrně jednoduché a mají podobu řádků oddělených čárkou. 
To umožňuje snadné algoritmické zpracování na straně počítače a zároveň jednoduché ruční testování v sériovém terminálu.

\section{Řízení zdroje měřicího signálu a rozsahů}

Zdroj měřicího signálu je založen na obvodu AD9837. 
Mikrokontrolér jej ovládá softwarově generovaným sériovým rozhraním na pinech PA5, PA6 a PA8. 
Při změně frekvence firmware vypočítá ladicí slovo podle hodinového kmitočtu 16 MHz a zapíše jej do frekvenčního registru DDS. 
Výstup lze také softwarově povolit nebo ponechat v resetu.

Amplituda výstupního signálu je nastavována digitálním potenciometrem MCP4561 přes sběrnici I\textsuperscript{2}C. 
Firmware zapisuje hodnotu 0 až 255 do jeho volatilního registru jezdce. 
Stejnosměrná složka zdroje je nastavována pomocí PWM výstupu na pinu PA7, který je v analogové části filtrován.

Proudový rozsah a zesílení obou měřicích větví jsou řízeny jako dvoubitové hodnoty. 
Piny PB8 a PB9 volí proudový rozsah, PB10 a PB11 zesílení proudového kanálu a PB12 s PB13 zesílení napěťového kanálu. 
Firmware pro tyto volby pouze nastavuje příslušné binární kombinace; fyzický význam jednotlivých kódů je dán zapojením analogové části.

\section{Synchronní snímání signálů}

Měřené střídavé signály jsou přivedeny na dva analogové vstupy mikrokontroléru. 
Napěťový kanál je připojen na PA1 a proudový kanál na PA4. 
Protože pro výpočet impedance je důležitý nejen poměr amplitud, ale i fázový rozdíl, musí být oba kanály vzorkovány současně.

Firmware k tomu využívá dvojici převodníků \acs{ADC} v režimu současné pravidelné konverze. 
ADC1 vzorkuje napěťový kanál a ADC2 proudový kanál. 
Spouštění převodu zajišťuje časovač TIM3, jehož událost update je přivedena jako externí trigger pro ADC1. 
Po každém impulzu časovače tedy vznikne jeden synchronní pár vzorků.

Výsledky převodu jsou ukládány pomocí \acs{DMA} do paměťového bufferu. 
Každá položka bufferu je 32bitové slovo, ve kterém je v dolních 16 bitech uložen vzorek napěťového kanálu a v horních 16 bitech vzorek proudového kanálu. 
Maximální délka jedné dávky je 2048 dvojic vzorků.

Vzorkování je dávkové. 
Hostitelský počítač pošle požadovaný počet vzorků a požadovanou vzorkovací frekvenci. 
Firmware podle požadované frekvence nastaví děličku a periodu časovače TIM3, spustí ADC a DMA, počká na dokončení přenosu a po skončení časovač i převodníky zastaví. 
Povolený rozsah požadované vzorkovací frekvence je 1 kHz až 1 MHz. 
Protože časovač používá celočíselné dělení, skutečná frekvence nemusí být přesně rovna požadované hodnotě; firmware proto v odpovědi vždy vrací skutečně nastavenou hodnotu.

Pokud je požadovaná vzorkovací frekvence v příkazu zadána jako nula, firmware ji odvodí z aktuální frekvence DDS. 
Cílem je přibližně 25 vzorků na periodu měřicího signálu, nejvýše však 1 MHz.

\section{Výstup naměřených dat}

Příkaz \texttt{capture} provede synchronní sejmutí obou kanálů, ale po sériové lince odešle pouze jeden z nich. 
Tento režim je vhodný hlavně pro rychlé zobrazení průběhu v jednoduchém osciloskopickém programu na počítači. 
Výstup má tvar
\begin{verbatim}
sample_rate_hz,200000
index,value
0,2048
1,2074
\end{verbatim}

Příkaz \texttt{capturepair} používá stejný snímací mechanismus, ale odešle oba kanály společně:
\begin{verbatim}
sample_rate_hz,200000
index,v_raw,i_raw
0,2048,2039
1,2074,2051
\end{verbatim}
Tento režim je důležitý pro další zpracování, protože zachovává dvojice vzorků ze stejného okamžiku.

\section{Výpočet impedance}

Příkaz \texttt{measure} provede sejmutí obou kanálů a následně v mikrokontroléru vypočítá amplitudu, fázi a komplexní impedanci. 
Nejprve jsou surové 12bitové hodnoty \acs{ADC} převedeny na napětí podle referenční hodnoty 3,3 V. 
Protože analogové signály jsou posunuty přibližně do poloviny napájecího rozsahu, odečítá se od nich hodnota 1,65 V.

Napěťový kanál je dále přepočten pomocí nastaveného zesílení napěťové větve. 
Proudový kanál je považován za napětí na zvoleném bočníku a po korekci zesílení proudové větve je převeden na proud. 
Ve firmwaru jsou použity hodnoty bočníků 100 $\Omega$, 1 k$\Omega$, 10 k$\Omega$ a 100 k$\Omega$. 
Tabulky zesílení napěťové a proudové větve jsou v aktuální verzi pouze zástupné a musí být nahrazeny hodnotami získanými měřením.

Amplituda a fáze jsou určeny výpočtem jedné složky diskrétní Fourierovy transformace na frekvenci nastaveného DDS. 
Pro každý kanál se počítají součty se sinem a kosinem
\begin{equation}
    X_c = \frac{2}{N} \sum_{n=0}^{N-1} x[n]\cos\left(2\pi f \frac{n}{f_s}\right),
\end{equation}
\begin{equation}
    X_s = \frac{2}{N} \sum_{n=0}^{N-1} x[n]\sin\left(2\pi f \frac{n}{f_s}\right),
\end{equation}
kde $N$ je počet vzorků, $f$ frekvence měřicího signálu a $f_s$ skutečná vzorkovací frekvence. 
Z těchto hodnot se určí amplituda
\begin{equation}
    A = \sqrt{X_c^2 + X_s^2}
\end{equation}
a fáze
\begin{equation}
    \varphi = \arctan2(-X_s, X_c).
\end{equation}

Fázový rozdíl je počítán jako rozdíl fáze napěťového a proudového kanálu. 
Protože je v aktuálním zapojení proudová analogová cesta invertovaná, firmware k výsledku přičítá kompenzaci 180$^\circ$. 
Komplexní impedance je následně určena z poměru amplitud a fázového rozdílu:
\begin{equation}
    |Z| = \frac{A_U}{A_I},
\end{equation}
\begin{equation}
    \operatorname{Re}(Z) = |Z|\cos(\varphi),
    \qquad
    \operatorname{Im}(Z) = |Z|\sin(\varphi).
\end{equation}

Výsledek příkazu \texttt{measure} obsahuje použitou vzorkovací frekvenci, frekvenci DDS, nastavený rozsah, amplitudu a fázi obou kanálů a vypočtenou reálnou i imaginární složku impedance.

\section{Stavová indikace}

Pro jednoduchou kontrolu běhu programu je použit výstup PB15. 
Firmware na něm generuje pomalý pulzující průběh, který slouží jako indikace, že hlavní smyčka běží a mikrokontrolér není zablokovaný.
